% =====================================================================
%  THÈME 5 — Rendement et bilan complet
%  Niveau FACILE — Corrigé
% =====================================================================
\documentclass[11pt,a4paper]{article}
\usepackage[utf8]{inputenc}
\usepackage[T1]{fontenc}
\usepackage{lmodern}
\usepackage[french]{babel}
\usepackage{amsmath,amssymb,mathtools}
\usepackage{icomma}
\usepackage[version=4]{mhchem}
\usepackage{siunitx}
\sisetup{output-decimal-marker={,},per-mode=symbol}
\usepackage{xcolor}
\usepackage{tikz}
\usepackage[most]{tcolorbox}
\usepackage{enumitem}
\usepackage{array}
\usepackage{fancyhdr}
\usepackage[a4paper,margin=2.1cm,headheight=26pt,headsep=8pt]{geometry}

\definecolor{primary}{RGB}{142,93,168}
\definecolor{primarydark}{RGB}{82,45,108}
\definecolor{corrcol}{RGB}{176,110,20}

\pagestyle{fancy}\fancyhf{}
\fancyhead[L]{\small\textcolor{primarydark}{\textbf{Fiche 5} — Thème 5 : Rendement}}
\fancyhead[R]{\small\textcolor{corrcol}{\textsc{Corrigé — Facile}}}
\fancyfoot[C]{\small\thepage}
\renewcommand{\headrulewidth}{0.8pt}
\renewcommand{\headrule}{\hbox to\headwidth{\color{primary}\leaders\hrule height \headrulewidth\hfill}}

\newtcolorbox[auto counter]{cor}[1][]{enhanced,breakable,boxrule=1pt,arc=3pt,
  left=3mm,right=3mm,top=2.4mm,bottom=2.4mm,colback=corrcol!6,colframe=corrcol,
  fonttitle=\bfseries,coltitle=white,
  title={Corrigé~\thetcbcounter\ \ifx\relax#1\relax\relax\else\textmd{--- \itshape #1}\fi}}

\newcommand{\rep}[1]{\par\smallskip\noindent\textbf{\color{corrcol}$\Rightarrow$ #1}}

\setlength{\parindent}{0pt}
\setlength{\parskip}{3pt}

\begin{document}

\begin{center}
{\LARGE\bfseries\color{primarydark}Thème 5 — Rendement}\\[1pt]
{\large\color{corrcol}\textbf{Corrigé — Niveau Facile}}\\
{\color{primary}\rule{0.6\linewidth}{1pt}}
\end{center}

\begin{cor}[rendement à partir de moles]
\[
R=\frac{n_{\text{réel}}}{n_{\text{théorique}}}\times100=\frac{0{,}080}{0{,}100}\times100=\boxed{\SI{80}{\percent}}.
\]
\end{cor}

\begin{cor}[rendement à partir de masses]
Comme $R=\dfrac{m_{\text{réel}}}{m_{\text{théorique}}}\times100$ (la masse molaire se simplifie) :
\[
R=\frac{36}{45}\times100=\boxed{\SI{80}{\percent}}.
\]
\end{cor}

\begin{cor}[quantité réelle à partir du rendement]
\[
n_{\text{réel}}=\frac{R}{100}\times n_{\text{théorique}}=\frac{75}{100}\times2{,}0=\boxed{\SI{1,5}{\mole}}.
\]
\end{cor}

\begin{cor}[masse réelle à partir du rendement]
\[
m_{\text{réel}}=\frac{R}{100}\times m_{\text{théorique}}=\frac{60}{100}\times50=\boxed{\SI{30}{\gram}}.
\]
\end{cor}

\begin{cor}[vrai ou faux ? notion de rendement]
\begin{enumerate}[label=\arabic*.,leftmargin=1.8em,topsep=1pt,itemsep=1pt]
  \item \textbf{Faux.} Par définition $n_{\text{réel}}\le n_{\text{théorique}}$, donc $R\le\SI{100}{\percent}$.
        On ne peut pas obtenir plus que le maximum stœchiométrique.
  \item \textbf{Faux.} L'excès permet seulement de consommer tout le limitant ; il n'agit pas sur les
        pertes ni les réactions secondaires. Le rendement peut rester faible.
  \item \textbf{Vrai.} $n_{\text{réel}}$ est une mesure expérimentale ; l'écart au théorique traduit
        les imperfections réelles de la réaction.
\end{enumerate}
\end{cor}

\end{document}
